\chapter{Background}
\label{chap:background}

\section{Motivation}
\label{sec:motivation}

An agent acting in its environment receives observational input, and responds to it with a commanded
control input.
Differently said, the agent is a function in relation to its environment, and to find an optimal
behavior, we need to approximate this optimal function as closely as possible. While Neural Networks
are known to be universal approximators~\cite{Cybenko1989Approximation,Hornik1989Multilayer}, and
hence able to approximate any function, in the case of an embodied agent, different constraints need to
be taken into account. Ay~\cite{Ay2015Geometric} names additionally to external embodiment constraints
possible intrinsic constraints, and argues that such can even be ``beneficial'' for the control
structure. Additionally, less parameters for deployment are helpful due to energy usage and storage
capacity, the core of the notion of cheap design\needcite{TODO}.
So, for such embodied agents, the goal is to find an ``embodied universal
approximator''~\cite[p.~3]{Montufar2015CheapControl}, which restricts the neural network
architecture to an intrinsic, unknown dimension, encoding physical properties belonging to the agent
and the environment, while disallowing impossible movements~\cite[Sec.~3]{Montufar2015CheapControl}.

Biology motivates this approach even further: Here it is clear to see, that by commanding a body,
some features of the control architecture can be simplified and delegated to the
embodiment\needcite{TODO}. For example, McGeer~\cite{McGeer1990Passive}'s passive walker clearly
shows, that some movements can be produced by the body without any control input, just by choosing
the right embodiment. While obviously here being simplified to not include any control at all, it
can quite simply be extended with such \needcite{\#9, reference to his previous paper}.
M{\"u}ller and Hoffmann~\cite[p.~4]{Mueller2017Morphological} mention such combination, as a robot
climbing up walls, whose design was inspired by geckos, indeed using energy for control, however
very little compared to an approach not exploiting these specific embodiment design. Further, it can
be seen that the complex
behaviors are not only less energy intensive to produce, but also utilize muscle
synergies~\cite{DAvella2003Combinations} and thus fewer control parameters. And while it is
difficult to clearly attribute the effect of the embodiment~\cite{Zahedi2013Quantifying}, for us,
this gives us the reason to believe that the simplicity, and thus efficiency, comes from a
low-dimensionality of the controller~\cite{Montufar2015CheapControl}.

That this is indeed true for certain examples, has been shown both in theory and practice
before~\cite{Datar2025Universal,Montufar2015Geometry,Li2018Measuring}. Li et
al.~\cite{Li2018Measuring} call this lower-dimension ``intrinsic dimension'', which is the
dimensionality required for the model to be able to find policies. They show for example that for achieving 90 percent of the performance of
the full model, the dimensionality can be reduced by 99.6 percent. For a convolutional neural
network, it is still 99.4 percent. Others show that the ``least complex model which is able to learn
an optimal policy does learn faster''~\cite{Montufar2015Geometry}. For a neural network reducing
parameters can be done in different ways\needcite{TODO}, some of which we will take a look at in
\cref{chap:inflation}. In comparison, for LQR, where the controller is a gain matrix, the
corresponding approach is reducing its rank~\cite{Recht2010Guaranteed}. However, while these
approaches are more parameter efficient while achieving similar, or slightly worse performance, it
was also observed that they tend to perform even worse when encountering
noise~\cite{Datar2025Universal}.

Robustness to noise however, is one of the important characteristics of controllers acting under
real-world conditions~\cite{Rao2001Naive,Roy2021MachineLearning}. One approach of achieving this
in nature is the principle of ``redundancy and diversity''~\cite{Kitano2004Biological}, which stands
in direct contrast to our goal of reducing redundant parameterizations. Another reason for producing
robust controllers is the difficult transfer from digital training environment to real-world
embodiments, commonly referred to as the ``reality gap''~\cite{Tobin2017Domain,Peng2018SimToReal}.
Approaches of training robust controllers typically introduce noise already during training, be it
either by using a LQG or including variations to the dynamics~\cite{Rao2001Naive,Peng2018SimToReal}.
These shown approaches however, don't look at low-dimensionality, which makes sense as
Kitano~\cite[p.~8]{Kitano2004Biological} poses that robustness automatically leads to ``increased resource
demands''.

Recent advancements in the research for low-rank policies for LQR\ldots

\section{Approach}
\label{sec:approach}

What I will do in the following sections

% reason why SVD is the right choice #23
% #17 bow tie structure in biological systems
